\hypertarget{appendix-l-100-more-interview-questions-part-5-8}{%
\section{Appendix L: 100 More Interview Questions (Part
5-8)}\label{appendix-l-100-more-interview-questions-part-5-8}}

These questions are designed to separate the ``Senior Engineers'' from
the ``Gods.'' If you can answer these without looking at the notes, you
are ready for any HFT or Systems Architecture interview on the planet.

\hypertarget{part-5-the-c-memory-model-atomics}{%
\subsection{Part 5: The C++ Memory Model \&
Atomics}\label{part-5-the-c-memory-model-atomics}}

\hypertarget{what-is-the-difference-between-stdmemory_order_relaxed-and-stdmemory_order_seq_cst}{%
\subsubsection{\texorpdfstring{1. What is the difference between
\texttt{std::memory\_order\_relaxed} and
\texttt{std::memory\_order\_seq\_cst}?}{1. What is the difference between std::memory\_order\_relaxed and std::memory\_order\_seq\_cst?}}\label{what-is-the-difference-between-stdmemory_order_relaxed-and-stdmemory_order_seq_cst}}

\textbf{Answer}: \passthrough{\lstinline!seq\_cst!} (Sequentially
Consistent) provides a global total ordering of all operations. It is
the safest but slowest. \passthrough{\lstinline!relaxed!} only
guarantees atomicity of the operation itselfit provides no guarantees
about the order of other memory operations.

\hypertarget{explain-release-acquire-semantics.}{%
\subsubsection{2. Explain ``Release-Acquire''
semantics.}\label{explain-release-acquire-semantics.}}

\textbf{Answer}: A \passthrough{\lstinline!memory\_order\_release!}
store ``synchronizes-with'' a
\passthrough{\lstinline!memory\_order\_acquire!} load of the same
variable. All memory writes performed by the storing thread
\emph{before} the release store are guaranteed to be visible to the
loading thread \emph{after} the acquire load.

\hypertarget{what-is-a-fences-memory-barrier}{%
\subsubsection{3. What is a ``Fences'' (Memory
Barrier)?}\label{what-is-a-fences-memory-barrier}}

\textbf{Answer}: A fence is an instruction that prevents the CPU or
compiler from reordering instructions across the fence boundary.
\passthrough{\lstinline!std::atomic\_thread\_fence!} can be used to
establish synchronization without a specific atomic variable.

\hypertarget{what-is-the-aba-problem-in-lock-free-programming}{%
\subsubsection{4. What is the ABA problem in lock-free
programming?}\label{what-is-the-aba-problem-in-lock-free-programming}}

\textbf{Answer}: It occurs when a thread reads a value A, another thread
changes it to B and then back to A. The first thread thinks nothing has
changed, but it might have (e.g., a node in a linked list was deleted
and a new one was allocated at the same address). \textbf{Fix}: Use
versioned pointers (hazard pointers) or
\passthrough{\lstinline!std::atomic<T>::compare\_exchange\_strong!} with
a counter.

\hypertarget{why-is-compare_exchange_weak-used-in-a-loop-instead-of-strong}{%
\subsubsection{\texorpdfstring{5. Why is
\texttt{compare\_exchange\_weak} used in a loop instead of
\texttt{strong}?}{5. Why is compare\_exchange\_weak used in a loop instead of strong?}}\label{why-is-compare_exchange_weak-used-in-a-loop-instead-of-strong}}

\textbf{Answer}: On some architectures (like ARM/Load-Link
Store-Conditional), \passthrough{\lstinline!weak!} can fail spuriously
even if the values match. However, \passthrough{\lstinline!weak!} is
faster in a loop because it allows the compiler to generate more
efficient code.

\begin{center}\rule{0.5\linewidth}{0.5pt}\end{center}

\hypertarget{part-6-lock-free-structures-concurrency}{%
\subsection{Part 6: Lock-Free Structures \&
Concurrency}\label{part-6-lock-free-structures-concurrency}}

\hypertarget{implement-a-lock-free-stack-treiber-stack.}{%
\subsubsection{6. Implement a Lock-Free Stack (Treiber
Stack).}\label{implement-a-lock-free-stack-treiber-stack.}}

\begin{lstlisting}[language={C++}]
template <typename T>
class LockFreeStack {
    struct Node { T data; Node* next; };
    std::atomic<Node*> head;
public:
    void push(T val) {
        Node* newNode = new Node{val, head.load()};
        while (!head.compare_exchange_weak(newNode->next, newNode));
    }
};
\end{lstlisting}

\hypertarget{what-is-false-sharing-and-how-do-you-prevent-it-in-c17}{%
\subsubsection{7. What is ``False Sharing'' and how do you prevent it in
C++17?}\label{what-is-false-sharing-and-how-do-you-prevent-it-in-c17}}

\textbf{Answer}: It happens when two independent atomic variables reside
on the same CPU cache line. Updating one invalidates the cache for the
other core. \textbf{Fix}: Use
\passthrough{\lstinline!alignas(hardware\_destructive\_interference\_size)!}
from \passthrough{\lstinline!<new>!}.

\hypertarget{explain-the-double-checked-locking-pattern-and-why-it-was-broken-before-c11.}{%
\subsubsection{8. Explain the ``Double-Checked Locking'' pattern and why
it was broken before
C++11.}\label{explain-the-double-checked-locking-pattern-and-why-it-was-broken-before-c11.}}

\textbf{Answer}: It was broken because the compiler could reorder the
object allocation and the pointer assignment, leading a second thread to
see a non-null pointer to an uninitialized object. C++11's memory model
(and \passthrough{\lstinline!std::atomic!}) fixed this.

\begin{center}\rule{0.5\linewidth}{0.5pt}\end{center}

\hypertarget{part-7-template-metaprogramming-tmp}{%
\subsection{Part 7: Template Metaprogramming
(TMP)}\label{part-7-template-metaprogramming-tmp}}

\hypertarget{what-is-sfinae-give-a-concrete-example.}{%
\subsubsection{9. What is SFINAE? Give a concrete
example.}\label{what-is-sfinae-give-a-concrete-example.}}

\textbf{Answer}: ``Substitution Failure Is Not An Error.'' It allows the
compiler to discard a template overload if the type substitution fails,
instead of throwing a hard error.

\begin{lstlisting}[language={C++}]
template <typename T>
auto func(T t) -> decltype(t.push_back(0)) { ... } // Only works for containers
\end{lstlisting}

\hypertarget{how-do-c20-concepts-improve-upon-sfinae}{%
\subsubsection{10. How do C++20 Concepts improve upon
SFINAE?}\label{how-do-c20-concepts-improve-upon-sfinae}}

\textbf{Answer}: Concepts provide a formal, readable way to constrain
templates. Instead of cryptic template vomit, you get clear errors:
``Type X does not satisfy requirement `HasPushBack'.''

\hypertarget{what-is-the-curiously-recurring-template-pattern-crtp}{%
\subsubsection{11. What is the Curiously Recurring Template Pattern
(CRTP)?}\label{what-is-the-curiously-recurring-template-pattern-crtp}}

\textbf{Answer}: A pattern where a class
\passthrough{\lstinline!Derived!} inherits from
\passthrough{\lstinline!Base<Derived>!}. It allows for ``Static
Polymorphism''achieving polymorphic behavior without the cost of virtual
functions.

\hypertarget{explain-stdvoid_t-and-how-its-used-for-trait-detection.}{%
\subsubsection{\texorpdfstring{12. Explain \texttt{std::void\_t} and how
it's used for trait
detection.}{12. Explain std::void\_t and how it's used for trait detection.}}\label{explain-stdvoid_t-and-how-its-used-for-trait-detection.}}

\textbf{Answer}: \passthrough{\lstinline!void\_t!} is a template that
always maps any list of types to \passthrough{\lstinline!void!}. It's
used to check if a certain member or type exists within a class during
template instantiation.

\begin{center}\rule{0.5\linewidth}{0.5pt}\end{center}

\hypertarget{part-8-systems-performance}{%
\subsection{Part 8: Systems \&
Performance}\label{part-8-systems-performance}}

\hypertarget{what-is-rtti-and-why-do-hft-developers-often-disable-it}{%
\subsubsection{13. What is RTTI and why do HFT developers often disable
it?}\label{what-is-rtti-and-why-do-hft-developers-often-disable-it}}

\textbf{Answer}: Runtime Type Information. It powers
\passthrough{\lstinline!dynamic\_cast!} and
\passthrough{\lstinline!typeid!}. It's disabled
(\passthrough{\lstinline!-fno-rtti!}) to save space in the binary and
avoid the overhead of storing type info in the vtable.

\hypertarget{what-is-the-difference-between-inline-and-__attribute__always_inline}{%
\subsubsection{\texorpdfstring{14. What is the difference between
\texttt{inline} and
\texttt{\_\_attribute\_\_((always\_inline))}?}{14. What is the difference between inline and \_\_attribute\_\_((always\_inline))?}}\label{what-is-the-difference-between-inline-and-__attribute__always_inline}}

\textbf{Answer}: \passthrough{\lstinline!inline!} is just a suggestion;
the compiler can ignore it. \passthrough{\lstinline!always\_inline!} (a
GCC/Clang intrinsic) forces the compiler to inline the function unless
it's physically impossible.

\hypertarget{explain-instruction-cache-warming.}{%
\subsubsection{15. Explain ``Instruction Cache
Warming.''}\label{explain-instruction-cache-warming.}}

\textbf{Answer}: It's the practice of running a piece of code (like a
trading strategy) with ``dummy data'' before the market opens, just to
ensure the instructions are loaded into the CPU's L1-Instruction cache.

\begin{center}\rule{0.5\linewidth}{0.5pt}\end{center}

\emph{Note: This is just the beginning. The next 85 questions in your
journey will cover everything from SIMD intrinsics to Linux Kernel
tuning. Keep pushing. The machine is waiting.}
